\section{Institutional background}\label{hist}
Public procurement in the state of São Paulo is conducted through a decentralized system of purchasing agencies, each referred to as a Public Buyer Unit (PBU). Since 2007, the use of the state's electronic procurement platform, the \emph{Bolsa Eletrônica de Compras} (BEC), has been mandatory for PBUs when acquiring common goods and services. The BEC platform centralizes tender procedures and records detailed information on each procurement process. It plays a crucial role in São Paulo's procurement operations: for instance, in 2018 alone, approximately US\$1.7~billion worth of goods and services were procured via BEC. Over the period from 2009 onward, the BEC system handled more than R\$7.4~billion (approximately 34\% of all state government purchases) for the State Health Secretariat (SES/SP) and other agencies in São Paulo. Given the large scale and value of these tenders, ensuring fairness and preventing collusion in the procurement process is of paramount importance. A failed tender (one that ends with no supplier selected) not only delays public service delivery but also wastes significant administrative resources, as each tender process can cost the government between US\$500 and US\$5,200 in preparation and execution.

The BEC catalog contains a comprehensive listing of standardized items and services, classified by category. Nearly 55,000 purchase offers were issued through BEC during our study period, covering about 6,350 distinct items ranging from medical supplies to office equipment. Each purchase offer (PO) is a digital solicitation document that identifies one or more items to be procured, along with the required quantities and specifications. Each item within a PO is auctioned separately; thus a unique procurement transaction can be defined by a combination of a PO number and an item code (POI). The BEC records include rich information for each POI, including the reference (reserve) price set by the buyer, the number of participating firms, all bids submitted (identifying winners and losers), the auction format used, and the identities of the buyer (PBU) and auctioneer.

São Paulo's PBUs predominantly use two types of competitive tendering procedures available on BEC: (i) a sealed-bid tender (\emph{convite}) and (ii) a multi-round descending auction (\emph{pregão}). In a sealed-bid tender, firms submit confidential bids by a specified deadline. After the submission period closes, all bids are opened and the contract is awarded to the qualified bidder with the lowest bid price. The sealed-bid format is typically used for purchases up to a value limit of R\$176,000 (approximately US\$35,000). In contrast, the descending auction (pregão) has no value limit and involves an open, interactive bidding phase. Initially, firms submit sealed proposals which are evaluated for qualification. Then, all qualified bids are announced (with the bidders' identities concealed) and an open auction phase commences. During this phase, which lasts for an initial 20 minutes, bidders can submit progressively lower bids, knowing the current lowest bid in real time. If a new bid is placed in the final minutes of this period (between 16 and 20 minutes), the auction is automatically extended by an additional 4 minutes. The auction concludes when no new bids arrive for 4 consecutive minutes. The winner is the bidder offering the lowest price at the end of this process, provided that price is below the government's preset reference price.

The key difference between the sealed-bid and descending auction formats lies in the opportunity for iterative bidding and post-auction negotiation. The pregão format, after determining the lowest bid through the auction, often allows a brief negotiation or confirmation step to potentially further reduce the price or clarify terms with the lowest bidder, as permitted by Brazilian procurement rules. In contrast, the convite (sealed bid) format is a single-shot process with no negotiation stage, which makes it simpler but also means the outcome is entirely determined at bid opening. From an oversight perspective, the open auction format provides more transparency during the bidding process (as all bids are observable in real time, albeit anonymously), whereas the sealed bid format concentrates all transparency at the moment of opening. Both mechanisms are subject to anti-collusion regulations and audit processes, but their procedural differences can influence how collusion or favoritism might manifest. Our analysis will focus specifically on sealed bid auctions. %subsequent sections will consider these institutional features when interpreting the results.

